\section{Executive summary}

\id{Summarize the assignment, how it went, general impressions, etc. Delete these blue comments from the final version of your report.}

\section{Requirements traceability}

\begin{table}[h]
    {\small
    \begin{tabular}{@{}p{2cm}p{2cm}p{2cm}p{8cm}@{}}
        \toprule
        \multicolumn{1}{c}{\textbf{Req ID}} & \multicolumn{1}{c}{\textbf{Status}} & \multicolumn{1}{c}{\textbf{Implemented in}} & \multicolumn{1}{c}{\textbf{Design considerations}} \\
        R1.1 & \{Missing, Partial, Implemented\} & \code{MapSetup.java} & The software sets up the same map by a hard-wired specification. The specification is given in the \code{MapSetup} class' \code{tiles} data structure and used in the \code{generateMap()} method at the initialization phase of the program. \\
        ... & ... & ...  & ... \\ \bottomrule
    \end{tabular}}
\end{table}

\id{The goal of this matrix is to aid your TA verify the satisfaction of requirements. If the table is superficial or too verbose, your TA will have a hard time verifying the completeness of your deliverable, and that will result in lost marks.
Use the Mising-Partial-Implemented status markers in your final report. Use the code{} command to format file names and Java code elements, such as class names, method names, etc.}

\section{Model a small domain in a drawing tool}

\textbf{Document the syntax of your model. This should be, basically, a list in which
enumerate every shape, arrow, etc and explain what they are. This is the concrete
syntax, i.e., what do we see when we look at the model. You should define the
abstract syntax, too, i.e., the underlying essential concepts and their relationships.}

\section{Propose an algorithm to translate models described in your syntax to Java or
Python code}

\textbf{Document the algorithm in a pseudocode and explain it briefly. Did you draw on a
previously studied content? If so, which course? If not, did you find the algorithm
somewhere or did you come up with it?}

\section{Implement a proof-of-concept solution}



\section{Reflections}

\textbf{Assess the complexity of extending your code generator to generate code in the other language: in Python, if you generated Java code in Tasks 2-3; and Java, if you generated Python code in Tasks 2-3. Would this be a complex task or a fairly straightforward one? How about generating C code? What is the main difference between extending your generator to Java or Python vs to C?}

Extending the code generator in Python instead of Java would be fairly simple. The overall structure of classes, inheritance, and fields is similar. The main change would be adjusting syntax, such as removing type declarations, changing constructors, and adapting list structures. The conceptual model to code would remain mostly the same. 
Generating Java or Python is easier than C. C does not support classes, inheritance, or objects directly so many concepts in this model would need to be simulated using structure, function pointer, and memory management. This makes the translation more complex and less direct. Java and Python are OOP languages, which map naturally to class diagrams, while C is primarily procedural, so the conceptual model does not translate so easily. 

\textbf{Code generators for common design formalisms (e.g., Class Diagrams to Java, Simulunk to C) are readily available. There are domains, however, in which new design formalisms pop up frequently (in order to support experts in expressing their designs in new ways). Such design languages and their supporting tooling are developed in so-called “language workbenches.” Research the Eclipse ecosystem to find the two top language workbenches. One of them allows you to define visual formalisms, the other one allows you to design textual formalisms. The latter are very similar to programming languages. In fact, they can integrate with them (not just generate code for them). Which are these two language workbenches and how would you use them in your term assignment (Catan)?}

Two major language workbenches in the Eclipse ecosystem are Sirius Web and Capella. 
Sirius Web is a web-based tool used to create visual modelling languages and domain specific languages (DSLs). It allows developers to design custom diagrams and modelling environments for specific systems. 
Capella is an open-source Model-Based Systems Engineering (MBSE) tool that supports system modelling using the Arcadia method. It provides a structured environment for designing complex systems. 
In Catan, these tools could be used to visually model components like players, tiles, buildings, and relationships before writing that Java code. They could help organize the system design and potentially generate parts of the code. This would make the design clearer, more structured and easier to extend if the game became more complex. 


\section{Roles and responsibilities}

The team members contributed equally to the deliverable. \id{This is the ideal situation, but in case the workload has not been equal, please, report accordingly.}

\begin{itemize}
    \item \TODO{\textbf{John Doe} contributed to the conceptual design and the implementation of RQ1.1.}
    \item \TODO{...}
\end{itemize}
